
\textbf{Positive}: Benchmark rotation infrastructure reduces research waste on saturated evaluations. Early warning mechanism (2--6 year lead time) enables proactive benchmark governance. Conference organizers can integrate saturation dashboards into submission systems, providing benchmark selection rationale disclosure.

\textbf{Negative}: Premature deprecation risk if thresholds are miscalibrated. Community coordination challenges exist---who decides when to rotate? Potential gaming: researchers might avoid ''risky'' saturated benchmarks even when scientifically appropriate.

\textbf{Mitigation}: Conservative thresholds reduce false positives. Expert consensus validation provides community input rather than purely algorithmic decisions. Voluntary adoption (no enforcement) allows researchers to override saturation signals when justified.
